II) \emph{Caso en que $\dim(A)=1$.} --- Sea $A'$ la clausura integral de
$A$; se sabe por (\textbf 0, 23.1.6) que $A'$ es una $A$-álgebra \emph{finita} y un
anillo \emph{local} completo. Si se pone $B'=B\otimes_A A'$, se tiene
$B\otimes_A K=B'\otimes_{A'}K$. El mismo razonamiento que en I) muestra
que, para todo ideal maximal $\mathfrak n'$ de $B'$, basta demostrar que
los anillos locales de $B'_{\mathfrak n'}\otimes_{A'}K$ satisfacen
$\boldsymbol R$; además, ese razonamiento muestra también que $A'$ y
$B'_{\mathfrak n'}$ satisfacen las hipótesis (i) y (ii) del enunciado y
que se cumple
$\boldsymbol P(\operatorname{Spec}(B'_{\mathfrak n'}\otimes_{A'}k'),k')$
(donde $k'$ es el cuerpo residual de $A'$). Además, como
$\dim(A')=1$ (\textbf 0, 16.1.5) y $A'$ es íntegro e integralmente cerrado,
es un \emph{anillo de valoración discreta completo}. Es posible, pues,
limitarse al caso en que $A$ ya es un anillo de valoración discreta
completo. Si $u$ es entonces un generador del ideal maximal de $A$,
el hecho de que $u$ sea $A$-regular y de que $B$ sea un $A$-módulo
plano implica que $t=\varphi(u)$ es un elemento $B$-regular que
pertenece al ideal maximal de $B$. La hipótesis de que se cumple
$\boldsymbol P(\operatorname{Spec}(B\otimes_A k),k)$ implica, en
particular, que se cumple $\boldsymbol R(B/tB)$; en virtud de
$(\mathrm R_{\mathrm{IV}})$, se cumple, por tanto, $\boldsymbol R(B)$.
En otras palabras, $U_{\mathbf R}(\operatorname{Spec}(B))$ contiene
el punto cerrado de $\operatorname{Spec}(B)$. Pero
$U_{\mathbf R}(\operatorname{Spec}(B))$ es abierto en virtud de
$(\mathrm R_{\mathrm{III}})$, y $\operatorname{Spec}(B)$ es el único
abierto de $\operatorname{Spec}(B)$ que contiene el punto cerrado.
Por ello, todos los anillos locales de $\operatorname{Spec}(B)$ poseen
la propiedad $\boldsymbol R$, en particular los de
$\operatorname{Spec}(B\otimes_A K)$.
